\section{Cancellation and the matrix-nil algebra}
\label{sec:cancellation}

We now use the hierarchy of Proposition~\ref{prop:hierarchy} to
construct the algebra needed for the group argument.  Fix its absolute
constants $c,k$, choose an integer
\begin{equation}
  D>10(c+1)(k+1),
  \label{ca:D}
\end{equation}
and then perform the program construction in that proposition and choose
$h_0$ sufficiently large.  Thus $D$ is fixed before the program fixed point
and before $h_0$.  Throughout this section put
\[
  h_j=h_0D^j,\qquad q_j=2^{h_j},\qquad
  L_j=2q_j+1,\qquad \Sigma_j=\Sigma_{h_j}.
\]
In particular $|\Sigma_j|\leq 2^{ch_j}$.  Write $A_j(K)=A_{h_j}(K)$ for
the algebra at level $j$ over a field $K$ of characteristic two, and write
\[
  E_j:A_{j+1}(K)\longrightarrow A_j(K)
\]
for the homomorphism induced by the block substitution $E_{h_j}$.  The
same symbol denotes that substitution on words.  Each upper letter is
replaced by a word of length $L_j$.

For a word $w$ on $\Sigma_j$, let $[w]_j$ denote its algebra image.  The
empty word has image $1$.  Scalar extension gives
$A_j(K)=K\otimes_{\F_2}A_j(\F_2)$: the defining relations have coefficients
in $\F_2$, and tensoring vector spaces over $\F_2$ is exact.  In particular
a nonzero word over $\F_2$ remains nonzero over $K$.  We use the usual
identification
\[
  M_d(K)\otimes_K A_j(K)=M_d(A_j(K)).
\]
Thus a numerical matrix commutes with an algebra element when the latter
is placed in every diagonal position, but numerical matrix factors retain
their original order.

\subsection{Weighted sums and matrix powers}

A matrix with entries in the positive-degree part of $A_j(K)$ has a
finite expansion
\[
 T=\sum_{\gamma\in\mathcal C}
 A_\gamma\otimes[v_\gamma]_j,
 \qquad A_\gamma\in M_d(K),\quad |v_\gamma|\geq1.
\]
Each term in $T^a$ is indexed by a sequence of $a$ summands, whereas its
total word length is the sum of the lengths of their labels. This length
can vary among terms in the same power. We first prove
cancellation after grouping such sequences by total word length. The
final subsection will recover cancellation at each sufficiently large
fixed number of factors.

\begin{definition}[Weighted tests]
\label{ca:tests}
A \emph{letter test} of dimension $d\geq1$ at level $j$ assigns a matrix
$M(a)\in M_d(K)$ to each $a\in\Sigma_j$.  For
$w=a_1\cdots a_n$ put $M(w)=M(a_1)\cdots M(a_n)$ and $M(\varnothing)=1_d$.
Its length-$n$ sum is
\begin{equation}
  S_{j,n}(M)=\sum_{w\in\Sigma_j^n} M(w)\otimes[w]_j.
  \label{ca:letter-sum}
\end{equation}
A \emph{chunk test} of dimension $d$ and label bound $b\geq1$ consists of
a finite set $\mathcal C$ of chunk identifiers, a nonempty label
$v_\gamma\in\Sigma_j^*$ with $|v_\gamma|\leq b$ for each
$\gamma\in\mathcal C$, and a weight $W_\gamma\in M_d(K)$.  Its length-$n$
sum is
\begin{equation}
  C_{j,n}(\mathcal C,W)=
  \sum_{\substack{r\geq0,\ \gamma_1,\ldots,\gamma_r\in\mathcal C\\
                   |v_{\gamma_1}|+\cdots+|v_{\gamma_r}|=n}}
  W_{\gamma_1}\cdots W_{\gamma_r}
       \otimes[v_{\gamma_1}\cdots v_{\gamma_r}]_j.
  \label{ca:chunk-sum}
\end{equation}
The empty sequence contributes $1_d\otimes1$ when $n=0$.
A test \emph{eventually vanishes} if its length-$n$ sum is zero for every
sufficiently large integer $n$.
\end{definition}

Different chunk identifiers may have the same label.  Moreover, different
sequences may spell the same word with different segmentations.  All such
sequences occur separately in Equation~\eqref{ca:chunk-sum}.  The sums are
finite because every label has positive length. For a letter test,
the length-$n$ sum is exactly a matrix power:
\[
 S_{j,n}(M)=\left(\sum_{a\in\Sigma_j}M(a)\otimes[a]_j\right)^n.
\]
For longer chunk labels, the length-$n$ sum instead includes sequences
with different numbers of factors. Unless a different bound is specified,
a chunk test below has label bound $k$.

\subsection{Cancellation inside one palette run}

The cancellation estimate uses an elementary alternating-multilinear
argument. Its dimension bound is sufficient for the hierarchy below.

\begin{lemma}[Alternating matrix sum]
\label{ca:alternating}
For a field $K$ of characteristic two and matrices
$X_1,\ldots,X_m\in M_d(K)$, put
\[
  \mathcal A_m(X_1,\ldots,X_m)
   =\sum_{\sigma\in S_m}X_{\sigma(1)}\cdots X_{\sigma(m)}.
\]
If $m>d^2$, then $\mathcal A_m(X_1,\ldots,X_m)=0$.
\end{lemma}

\begin{proof}
The map $\mathcal A_m$ is multilinear.  If $X_i=X_j$ with $i\ne j$,
pair a permutation $\sigma$ with $(i\ j)\circ\sigma$.  The two products
are equal, and the pairing has no fixed point.  Their sum is zero in
characteristic two.  Hence $\mathcal A_m$ is alternating, meaning that
it vanishes whenever two arguments agree.  Expand its arguments in a
basis of the $d^2$-dimensional vector space $M_d(K)$.  If $m>d^2$, each
term in the multilinear expansion has a repeated basis argument and is
therefore zero.  This proves the result without any division by $m!$.
\end{proof}

\begin{lemma}[Direct cancellation]
\label{ca:direct}
Fix $j,d,b$.  If
\begin{equation}
  \frac{q_j}{100}>b(d^2+3),
  \label{ca:direct-bound}
\end{equation}
then every dimension-$d$ chunk test of label bound $b$ at level $j$
eventually vanishes over every field of characteristic two.
\end{lemma}

\begin{proof}
Fix such a test and choose $n$ above the long-word threshold in
Proposition~\ref{prop:hierarchy}.  In its length-$n$ sum discard the
sequences whose concatenated word is zero.  We partition the remaining
sequences into permutation orbits and prove that each orbit contributes
zero.

For each remaining sequence let $w$ be its concatenated word.  Erase only
the color field from each letter of $w$, retaining its position and all
other fields. Among the maximal eligible runs, choose the first one
having length at least $q_j/100$. This choice exists by
Proposition~\ref{prop:hierarchy} and depends only on the color-erased
word.  Every eligible run in a nonzero word has distinct colors: if two
colors agreed, adjacent transpositions within that run could bring them
together and give a declared zero substring.  This also explains why
color distinctness need not enter the rule selecting the run.

Mark the chunk boundaries of the sequence.  Its \emph{packet} is the
consecutive list of all complete chunks contained in the chosen run.
Keep all other chunks fixed, including any chunk crossing an endpoint of
the run.  There are at most two crossing chunks.  If the run has length
$r$ and the packet contains $m$ chunks, then
\[
  r\leq mb+2b,\qquad
  m\geq r/b-2>d^2.
\]
Here the last inequality follows from Equation~\eqref{ca:direct-bound}.
In particular the packet is nonempty.

Permute these $m$ whole chunks arbitrarily.  They may have unequal
lengths, but their total length is fixed.  Thus the packet occupies the
same letter interval after every permutation, and every letter in that
interval retains the common non-color tuple.  The color-erased word is
unchanged at every position.  Consequently the chosen maximal run is
unchanged.  The exterior chunks, including the endpoint-crossing chunks,
are unchanged as well, so the complete chunks in this run are again
exactly the packet just permuted.  Its internal chunk boundaries may
move; its endpoints and the rule selecting it do not.

All words obtained in this way have the same algebra image: arbitrary
permutations of their letters within the run are products of the allowed
adjacent transpositions.  If any permuted word were zero, their common
image would be zero and the starting word would also be zero.  Thus the
discarded zero sequences cannot remove any member of a nonzero orbit,
and every permuted sequence is still among the summands under
consideration.  Conversely applying
the same procedure to any of them produces the same set of sequences.
These sets are therefore a partition into orbits.

The labels of the packet chunks are pairwise distinct.  Indeed each
begins with a different color, since their first letters occupy distinct
positions of the distinct-color run.  Hence the chunk identifiers in the
packet are also pairwise distinct, even if the original list admitted
several identifiers with a common label.  All $m!$ permutations give
different sequences of identifiers; the orbit is a free permutation
orbit.  This assertion concerns sequences, not their concatenated words.
Distinct segmentations of a word remain separate summands and belong to
their respective orbits.

Choose one orbit and write its fixed exterior matrix products as $U,V$,
the packet weights as $X_1,\ldots,X_m$, and its common algebra image as
$a\in A_j(K)$.  Its entire contribution is
\[
  U\left(\sum_{\sigma\in S_m}
      X_{\sigma(1)}\cdots X_{\sigma(m)}\right)V\otimes a.
\]
The weights $X_i$ themselves need not be distinct.  The displayed middle
sum vanishes by Lemma~\ref{ca:alternating}.  There are only finitely many
sequences of total length $n$, so summing the zero orbit contributions
proves the desired vanishing.
\end{proof}

\subsection{Moving a test to a larger palette}

Direct cancellation applies when the palette is large compared with the
test dimension. The fixed palette at level zero cannot meet
Equation~\eqref{ca:direct-bound} for every dimension. We first treat a
letter test by moving it to the next level, where the palette is larger.
The cuts turn
it into a chunk test on the bounded upper words attached to tokens. The
matrix automaton then turns such a chunk test back into a letter test, so
that the passage to the next level can be repeated. We keep track of the
increase in dimension and apply direct cancellation as soon as its
inequality holds for an incoming chunk test. The automaton will also
reduce any level-zero chunk test to a letter test, thereby covering the
longer labels in the original matrix expansion.

\begin{lemma}[Reduction by the cuts]
\label{ca:cut-reduction}
Fix a level $j$, a field $K$ of characteristic two, and a dimension $d$.
If every dimension-$d$ chunk test of label bound $k$ at level $j+1$
eventually vanishes, then every dimension-$d$ letter test at level $j$
eventually vanishes.
\end{lemma}

\begin{proof}
Let $M$ be a letter test at level $j$.  By
Proposition~\ref{prop:hierarchy}, there are integers $B,N_0$ such that
every nonzero word of length at least $N_0$ has the stated cut
decomposition, with initial and final pieces of length at most $B$.
Increase $B$ to at least one.  A cut is recognized solely from the letter
immediately after it; there is no additional cut after the last letter of
the word.

Take the token list $\mathcal T$ to be a set of distinct literal words.
For each $t\in\mathcal T$ fix once and for all an upper word $u_t$ with
\begin{equation}
  1\leq |u_t|\leq k,\qquad
  |t|=L_j|u_t|,\qquad [t]_j=E_j([u_t]_{j+1}).
  \label{ca:token-normalization}
\end{equation}
If a token has more than one possible upper normal form, choose one;
it is the literal token, not its normal form, that indexes this list.
Every token begins with a cut letter and contains no further cut letter.

Let $\mathcal P$ be the finite set of words of length at most $B$ containing
no cut letter, including the empty word.  Let $\mathcal S$ be the finite
set of words of lengths $1$ through $B$ whose first letter is a cut letter
and whose remaining letters are not.  Consider concatenations
\begin{equation}
  p\,t_1\cdots t_r\,s,
  \qquad p\in\mathcal P,\quad s\in\mathcal S,\quad t_i\in\mathcal T.
  \label{ca:cut-language}
\end{equation}
They have a unique such representation: the prefix ends at the first cut,
the suffix begins at the last cut, and the intervening cut intervals are
exactly $t_1,\ldots,t_r$.  This also proves uniqueness when $r=0$.
For sufficiently large length every nonzero word has this representation.
Some concatenations in Equation~\eqref{ca:cut-language} may be zero, and
some zero words may have no representation; in either case their
contribution to Equation~\eqref{ca:letter-sum} is zero.  Thus summing all
concatenations in Equation~\eqref{ca:cut-language} neither loses nor
duplicates any nonzero contribution.

Use $\mathcal T$ as the chunk identifiers at level $j+1$, with label
$u_t$ and weight $M(t)$.  Denote the resulting length-$m$ sum by $C_m$.
For every sufficiently large $n$ we have the exact identity
\begin{equation}
  S_{j,n}(M)=
  \sum_{\substack{p\in\mathcal P,\ s\in\mathcal S\\
       m=(n-|p|-|s|)/L_j\in\Z_{\geq0}}}
  \bigl(M(p)\otimes[p]_j\bigr)
  (\mathrm{id}\otimes E_j)(C_m)
  \bigl(M(s)\otimes[s]_j\bigr).
  \label{ca:cut-identity}
\end{equation}
Indeed the length equation follows from
Equation~\eqref{ca:token-normalization}; multiplicativity of $M$ preserves
the matrix order, and the token equalities give the algebra factor.
No nonzero-preservation or injectivity assertion about $E_j$ is required
for this identity.  In particular, an upper concatenation whose encoded
product vanishes causes no difficulty.

By hypothesis $C_m=0$ for $m\geq N_1$, for some $N_1$.  Since
$|p|+|s|\leq2B$, every integer $m$ occurring in
Equation~\eqref{ca:cut-identity} is at least $N_1$ once
$n\geq2B+L_jN_1$.  Taking $n$ also large enough for that identity proves
the lemma.
\end{proof}

Finite-dimensional matrix products as word weights are part of the
theory of weighted automata and recognizable series developed by
Sch\"utzenberger \cite{Schutzenberger1961}. A related matrix encoding of
products of bounded lengths is used by Smoktunowicz
\cite[Sections~6 and~7]{Smoktunowicz2008} to derive quasi-regularity from
matrix nilpotence. The following prefix-state
construction keeps chunk boundaries explicit and supplies the dimension
estimate needed for our reduction.

\begin{lemma}[A matrix automaton for chunks]
\label{ca:automaton}
Let $\Sigma$ be a nonempty finite alphabet of size $s$, and let a chunk
test over $\Sigma$ have dimension $d$ and label bound $b$.  There is a
letter test of dimension at most
\begin{equation}
  C_b s^b d,\qquad C_b=b,
  \label{ca:automaton-bound}
\end{equation}
whose length-$n$ sum has the chunk sum as one fixed $d$ by $d$ block for
every $n$.  Consequently eventual vanishing of all letter tests in that
dimension implies eventual vanishing of the chunk test.
\end{lemma}

\begin{proof}
First combine all identifiers with the same label $v$ into one weight
\[
  W(v)=\sum_{\gamma:\,v_\gamma=v}W_\gamma.
\]
Distributivity shows that this leaves every chunk sum unchanged, including
the contributions from repeated occurrences of such a label.  Zero
weights can be retained or removed.  If there are no labels, take every
letter matrix to be the zero $d$ by $d$ matrix; this gives the assertion,
including length zero.  Hence suppose the label set is
$\mathcal V\ne\varnothing$.

The states are the empty word $\epsilon$ and all nonempty proper prefixes
of labels in $\mathcal V$.  For each state $p$ and letter $a\in\Sigma$
put an identity-weighted transition
\[
  p\xrightarrow{\ a,\,1_d\ }pa
\]
whenever $pa$ is a proper prefix of some label.  Put a closing transition
\[
  p\xrightarrow{\ a,\,W(pa)\ }\epsilon
\]
whenever $pa\in\mathcal V$.  Both kinds of transition are allowed when
$pa$ is both a label and a proper prefix.  We add the weights of any
transitions with identical initial state, final state, and letter.

If the number of states is $e$, define $B(a)\in M_{ed}(K)$ by taking its
$(p,q)$ block to be the sum of the weights of transitions from $p$ to $q$
with letter $a$.  The $(\epsilon,\epsilon)$ block of
$B(a_1)\cdots B(a_n)$ is the sum of the weights of all paths reading
$a_1\cdots a_n$ and beginning and ending at $\epsilon$.  The visits to
$\epsilon$ mark exactly the closures of chunks.  A path between two
consecutive visits reads one label $v$, contributes identity weights
until its last step, and then contributes $W(v)$.  Paths therefore
correspond bijectively to segmentations into labels, with exactly the
ordered matrix products required in Equation~\eqref{ca:chunk-sum}.
Taking the $(\epsilon,\epsilon)$ block after summing the words with their
algebra images proves the asserted identity.

Every state is a word of length less than $b$, so
\[
  e\leq\sum_{i=0}^{b-1}s^i\leq b s^{b-1}\leq C_b s^b.
\]
Padding the matrices by zero blocks, if needed, realizes any larger
specified dimension without changing positive-length sums in the original
block.  This proves the dimension claim.
\end{proof}

\subsection{Finite termination of the reduction}

A fixed level directly cancels tests whose dimension is small compared
with its palette. Passing to the next level enlarges the palette and
also enlarges the test dimension through the chunk automaton. We now
compare these two growth rates and stop the reduction after finitely
many levels for each fixed test.

\begin{proposition}[Eventual vanishing of every fixed test]
\label{ca:all-tests}
Over every field $K$ of characteristic two, every letter test at level
zero eventually vanishes.  The same is true of every level-zero chunk
test with any fixed finite label bound.
\end{proposition}

\begin{proof}
Fix the field and a starting letter-test dimension $d$.  For $j\geq1$
define integer dimension bounds
\begin{equation}
  d_j=d\prod_{i=1}^{j-1}\bigl(C_k|\Sigma_i|^k\bigr).
  \label{ca:dimension-product}
\end{equation}
The empty product gives $d_1=d$.  Lemma~\ref{ca:cut-reduction} reduces
dimension-$d$ letter tests at level zero to dimension-$d_1$ chunk tests
at level one.  At a level $i\geq1$, Lemma~\ref{ca:automaton} turns a
dimension-$d_i$ chunk test into a letter test of dimension at most
$d_{i+1}=C_k|\Sigma_i|^kd_i$ at that same level.  Padding to dimension
$d_{i+1}$ and applying Lemma~\ref{ca:cut-reduction} reduces its vanishing
to that of a dimension-$d_{i+1}$ chunk test at level $i+1$.

We show that these reductions can be stopped after finitely many levels.
Using $|\Sigma_i|\leq2^{ch_i}$ and the geometric sum for $h_i$, we obtain
\begin{align}
  \log_2(d_j^2)
  &\leq 2\log_2 d+2(j-1)\log_2 C_k
       +2kc\sum_{i=1}^{j-1}h_i \notag\\
  &=2\log_2 d+2(j-1)\log_2 C_k
       +\frac{2kc}{D-1}(h_j-h_1).
  \label{ca:dimension-growth}
\end{align}
By Equation~\eqref{ca:D}, the number
$\delta=1-2kc/(D-1)$ is strictly positive.  Since $h_j=h_0D^j$, it follows
from Equation~\eqref{ca:dimension-growth} that
\[
  h_j-\log_2(d_j^2)\longrightarrow+\infty.
\]
In particular there is a finite $J\geq1$ such that
\[
  2^{h_J}>400k d_J^2\geq100k(d_J^2+3),
\]
where the second inequality uses $d_J\geq1$.  Direct cancellation,
Lemma~\ref{ca:direct}, therefore applies to all dimension-$d_J$ chunk
tests of label bound $k$ at level $J$.

Work backwards through the finitely many reductions.  If all
dimension-$d_{i+1}$ chunk tests at level $i+1$ eventually vanish, then all
dimension-$d_{i+1}$ letter tests at level $i$ eventually vanish by
Lemma~\ref{ca:cut-reduction}.  Lemma~\ref{ca:automaton} then gives
eventual vanishing for all dimension-$d_i$ chunk tests at level $i$.
Descending to level one and using Lemma~\ref{ca:cut-reduction} once more
proves the letter-test assertion at level zero.

For precision, no common length threshold for all tests is being
postulated.  The stopping level $J$ above depends only on the starting
dimension and the fixed hierarchy.  Each application of the reductions
uses a finite test and, in the cut reduction, finitely many bounded end
pieces.  A threshold for the particular resulting test is pulled back
through finitely many length dilations and additions of bounded end
lengths.  The conclusion has the order of quantifiers
\[
  \text{for every fixed test, there is }N
  \text{ such that every length }n\geq N\text{ has zero sum}.
\]
It does not require one $N$ for all dimensions.

Finally, a level-zero chunk test of dimension $d$ with any label bound
$b$ is a block of a level-zero letter test of finite dimension at most
$C_b|\Sigma_0|^bd$, by Lemma~\ref{ca:automaton}.  The already proved
letter-test assertion applies to that test and proves the final claim.
\end{proof}

\subsection{The algebra and coefficient extraction}

\begin{theorem}
\label{thm:algebra}
There is an infinite, finitely presented, graded unital associative
$\F_2$-algebra
\[
  R=\bigoplus_{n\geq0}R_n=\F_2\oplus I,
  \qquad R_0=\F_2,\qquad I=\bigoplus_{n\geq1}R_n,
\]
such that $M_\ell(I)$ is nil for every integer $\ell\geq1$: each matrix
with entries in $I$ is nilpotent, with its own nilpotence index.
\end{theorem}

\begin{proof}
Take $R=A_0(\F_2)$, using only the level-zero alphabet and the level-zero
relations.  The alphabet is finite, and there are only finitely many
words of length at most $k$ on it.  The chosen replacement rules and
zero rules thus give an ordinary finite algebra presentation
\[
  R=\F_2\langle\Sigma_0\rangle/\mathcal J_0,
\]
where $\mathcal J_0$ is the two-sided ideal generated by $u-v$ for the
listed replacements $u\leftrightarrow v$ and by $w$ for the listed zero
words $w$.  No relation from a higher level is added to this ideal.  All
defining polynomials are homogeneous of positive degree, so word length
gives the displayed grading, the empty word spans $R_0=\F_2$, and the
positive-degree part $I$ is a two-sided ideal.

For every $j\geq1$, choose a nonzero letter $a_j$ at level $j$.  The word
\[
  w_j=E_0E_1\cdots E_{j-1}(a_j)
\]
is nonzero by repeated use of the nonzero-preservation clause of
Proposition~\ref{prop:hierarchy}.  Its length is
$L_0L_1\cdots L_{j-1}$, which tends strictly to infinity because each
$L_i>1$.  The nonzero classes $[w_j]_0$ lie in distinct homogeneous
degrees and are therefore linearly independent.  Consequently $R$ is
infinite-dimensional over $\F_2$, and in particular is an infinite set.

Fix $\ell\geq1$ and $T\in M_\ell(I)$.  Express its entries as finite
linear combinations of nonempty word images and collect matrix
coefficients to write
\begin{equation}
  T=\sum_{\gamma\in\mathcal C}
       A_\gamma\otimes[v_\gamma]_0,
  \qquad A_\gamma\in M_\ell(\F_2),\qquad |v_\gamma|\geq1.
  \label{ca:T-expansion}
\end{equation}
If this list is empty then $T=0$ and there is nothing to prove.
Otherwise let $b=\max_\gamma|v_\gamma|$.  Introduce a central
indeterminate $z$ and apply Proposition~\ref{ca:all-tests} over the field
$K=\F_2(z)$ to the chunks $v_\gamma$ with weights $zA_\gamma$.
Its length-$n$ sum is the image of the polynomial
\begin{equation}
  P_n(z)=
  \sum_{\substack{r\geq0,\ \gamma_1,\ldots,\gamma_r\in\mathcal C\\
                   \sum_i|v_{\gamma_i}|=n}}
    z^r A_{\gamma_1}\cdots A_{\gamma_r}
         \otimes[v_{\gamma_1}\cdots v_{\gamma_r}]_0
     \ \in M_\ell(R)[z].
  \label{ca:length-polynomial}
\end{equation}
For fixed $n$ this is a finite polynomial: nonempty chunks imply $r\leq n$
for every positive-length term.  The proposition supplies $N$ such that
the image of $P_n(z)$ in $M_\ell(K\otimes_{\F_2}R)$ is zero whenever
$n\geq N$.

The map
\[
  M_\ell(R)[z]\longrightarrow M_\ell(\F_2(z)\otimes_{\F_2}R)
\]
is injective.  For example, choose an $\F_2$-basis of $R$ and expand the
finitely many coefficients of a polynomial in that basis; its image is
zero exactly when every resulting scalar polynomial is zero in
$\F_2(z)$, hence in $\F_2[z]$.  Thus $P_n(z)$ itself is zero for all
$n\geq N$, and each of its coefficients can legitimately be extracted
in the original matrix algebra.

For an integer $a\geq1$, expansion of Equation~\eqref{ca:T-expansion}
and grouping its sequences by total word length gives the finite identity
\begin{equation}
  T^a=\sum_{n=a}^{ab}[z^a]P_n(z).
  \label{ca:power-coefficient}
\end{equation}
The lower bound $a$ uses that every factor has a nonempty label; the
upper bound $ab$ uses the fixed maximum label length.  If $a\geq N$,
every polynomial on the right of Equation~\eqref{ca:power-coefficient}
is zero.  Hence $T^a=0$.  Since $\ell$ and $T$ were arbitrary, every
$M_\ell(I)$ is nil, as asserted.
\end{proof}

The positive-degree presentation also gives a finitely presented algebra
without an identity. Recall that a nonunital algebra is nil if each of its
elements is nilpotent, whereas it is nilpotent if one power of the whole
algebra is zero. Its unitization adjoins a scalar identity; an algebra over
a field is algebraic if every element satisfies a nonzero polynomial over
that field.

\begin{corollary}[Finitely presented nil and algebraic algebras]
\label{cor:nil-algebra}
For the algebra $R=\F_2\oplus I$ of Theorem~\ref{thm:algebra}, the
following hold.
\begin{enumerate}
\item The positive-degree ideal $I$, regarded as a nonunital associative
$\F_2$-algebra, is finitely presented, infinite-dimensional, and nil, but
$I^m\ne0$ for every integer $m\geq1$. In particular, it is a finitely
presented Jacobson radical algebra that is not nilpotent.
\item The algebra $R$ is the unitization of $I$ and is an
infinite-dimensional finitely presented algebraic unital associative
$\F_2$-algebra. More precisely, if $r=a1+i$ with $a\in\F_2$ and $i\in I$,
there is an integer $N=N(i)\geq1$ such that
\[
 (r-a1)^N=i^N=0.
\]
Thus $r$ satisfies the nonzero monic polynomial $(t-a)^N$.
\end{enumerate}
\end{corollary}

\begin{proof}
Let $F=\F_2\langle\Sigma_0\rangle$ be the free unital algebra used in
Theorem~\ref{thm:algebra}, and let $F_+$ be its subalgebra spanned by the
nonempty words. This is the free nonunital associative algebra on the finite
alphabet $\Sigma_0$. Let $S\subset F_+$ be the finite set of defining
polynomials: the listed $u-v$ and the listed zero words $w$. Write $(S)_+$
for the two-sided ideal of the nonunital algebra $F_+$ generated by $S$.
Both $(S)_+$ and the unital ideal $\mathcal J_0=(S)_F$ are the linear span
of the products $usv$, with $s\in S$ and with each word context $u,v$
allowed to be empty. Indeed, closure under multiplication by $F_+$ gives
the nonempty contexts, while the generators themselves and one-sided
products give the empty contexts. Since every $s$ has positive degree,
these products lie in $F_+$. Hence $\mathcal J_0=(S)_+\subset F_+$, and
restriction of the quotient map gives
\[
 I\cong F_+/(S)_+.
\]
This is a finite presentation in the nonunital category. The splitting
$F=\F_2\oplus F_+$ descends to $R=\F_2\oplus I$ with the usual
unitization multiplication.

Theorem~\ref{thm:algebra} and this direct sum show that $I$ is
infinite-dimensional. Taking matrix size one in that theorem shows that
$I$ is nil. If $i^N=0$, then in its unitization
\[
 (1-i)^{-1}=\sum_{j=0}^{N-1}i^j.
\]
Thus every element of $I$ is quasi-regular, so $I$ is Jacobson radical.
If $I^m=0$ for some $m$, every word in the finite generating alphabet of
length at least $m$ would vanish. The images of the finitely many words
of lengths $1,\ldots,m-1$ would then span $I$, contradicting its infinite
dimension. Therefore $I^m\ne0$ for every $m\geq1$.

The finite presentation and infinite dimension of $R$ are already part of
Theorem~\ref{thm:algebra}. For $r=a1+i$ choose $N=N(i)$ with $i^N=0$.
Then $(r-a1)^N=i^N=0$, which proves that $r$ is algebraic over $\F_2$.
This applies to every element of $R$.
\end{proof}

\begin{remark}
\label{ca:no-uniform-index}
The nilpotence indices in Theorem~\ref{thm:algebra} and the polynomial
degrees in Corollary~\ref{cor:nil-algebra} may depend on the chosen matrix
or element. No common exponent is asserted.
\end{remark}
